\begin{textsample}{POS Dim 3 – pe – Score 13.00 – pe/pe\_inf\_23.txt}  \label{ex:f3_pos_186}
2.11 CAMPOS DE EXPERIÊNCIAS E DIREITOS DE APRENDIZAGEM E DESENVOLVIMENTO Conforme a BNCC, os campos de experiências “ constituem um arranjo curricular que \textbf{acolhe} as situações e as experiências concretas da vida cotidiana das \textbf{crianças} e seus saberes, entrelaçando -os aos conhecimentos que fazem parte do patrimônio cultural ” ( BRASIL, 2017, p. 38 ).
Essa abordagem tem foco nas experiências que são sustentadas pela lógica do conhecimento não como algo estanque, linear, instrumental, mas sim como ciência com uma construção que encanta e conecta a \textbf{criança} com o mundo, que favorece a aprendizagem e o desenvolvimento, que agrega diferentes linguagens, oportunizando possibilidades às \textbf{crianças} para explorarem, \textbf{brincarem}, participarem, conviverem, expressarem -se e conhecerem -se a si e ao mundo.
São cinco os campos de experiências denominados na BNCC ( 2017 ) : - O EU, O OUTRO E O NÓS ; - CORPO, \textbf{GESTOS} E MOVIMENTOS ; - TRAÇOS, \textbf{SONS}, CORES E FORMAS ; - \textbf{ESCUTA}, FALA, PENSAMENTO E IMAGINAÇÃO ; - ESPAÇOS, TEMPOS, QUANTIDADES, RELAÇÕES E TRANSFORMAÇÕES.
Essa forma de conceber o currículo na Educação \textbf{Infantil} se diferencia de uma organização curricular por áreas de conhecimento ou disciplinas escolares e tem respaldo na concepção dos Campos de Experiência, que são organizados a partir da articulação do conhecimento, das práticas sociais e das linguagens.
Conforme Fochi ( 2015 ), esses Campos apresentam : - As experiências concretas da vida cotidiana em virtude de nelas residirem situações importantes a serem consideradas e problematizadas para as \textbf{crianças}, tais como as atividades de \textbf{higiene}, alimentação, sono. - O convívio no espaço da vida coletiva nas interações com outras \textbf{crianças} e com os \textbf{adultos}. - A aprendizagem da cultura na articulação dos saberes das \textbf{crianças} com aqueles que a humanidade já sistematizou, na apropriação de rituais e modos de funcionamento de cada cultura. - A produção de narrativas individuais e coletivas através de diferentes linguagens, já que as \textbf{crianças} aprendem porque querem compreender o mundo em que vivem. - As \textbf{crianças} vivem suas \textbf{brincadeiras} de modo narrativo porque formulam e \textbf{contam} histórias ao mesmo tempo em que \textbf{dramatizam}.
É no contexto de cada Campo de Experiência que os direitos³ de aprendizagem e desenvolvimento ( conviver, \textbf{brincar}, participar, explorar, expressar e conhecer -se ) assumem seus respectivos significados e fundamentam os direitos e objetivos do currículo que se apresenta. ³ Os direitos de aprendizagem e desenvolvimento, que seguem os distintos Campos de Experiências, neste documento, são extraídos da 2ª versão da BNCC – 2016.

% matched lemmas: acolher, adulto, brincadeira, brincar, contar, criança, dramatizar, escuta, gesto, higiene, infantil, som
\end{textsample}
